\subsection{Chapter Officers} \label{I.1}

The officers of this chapter are President, Vice President,
Corresponding Secretary/Recording Secretary (hereafter referred to as Secretary),
Treasurer, two External Vice Presidents, Service Coordinator, Professional Development
Officer, Activities Officer, Graduate Student Vice President, Chapter Development Officer,
two K-12 Outreach Officers, Campus Outreach Officer, Membership Officer, Historian,
Publicity Officer, and all ad hoc officers listed in Appendix \ref{F.4}. All officers must be members.
In addition to the duties specified in the \href{http://www.tbp.org/off/ConstBylaw.pdf}{Tau Beta Pi Association Bylaws 5.03}, the
duties described in Appendix \ref{F} are required of the officers.


\subsection{Advisors} \label{I.2}

The chapter must maintain a minimum of four advisors, as stipulated
in \href{http://www.tbp.org/off/ConstBylaw.pdf}{C-VI,7} of the Constitution of the Tau Beta Pi Association. An advisor must be an
initiated member of Tau Beta Pi. Additionally, advisors must be a faculty member, a national
official, an alumni of MI-$\Gamma$ with appropriate chapter experience, a graduate student
with appropriate chapter experience, or a former chapter president. One of the advisors
must serve as the chapter’s Chief Advisor. This advisor will be selected by the advisors
biannually for a two year term. The advisors are responsible for ensuring chapter continuity.


\subsection{Advisory Board} \label{I.3}

The Advisory Board is composed of the Executive Committee
and the Chapter Advisors. The voting Executive Committee officers for the Advisory
Board are the President, Vice President, and Secretary. The voting advisors for the
Advisory Board are the Chief Advisor and the three most senior advisors.
\begin{enumerate}
    \item Seniority is measured as time spent in the current advisor term.
    \item In the event of a tie affecting determination of voting advisors, the President will
determine the voting advisor(s) from among those tied by a suitably random method
before the first contested vote of the semester.
\end{enumerate}


\subsection{Officer Corps} \label{I.4}

\begin{enumerate}
    \item \subsecitem{a}{Membership and Responsibilities}
    The Officer Corps consists of the officers and
advisors of the Chapter, and any other persons deemed necessary by the officers. The
officer corps:
    \begin{enumerate}
        \item \paritem{a.i}{Initiation Dues}
        Sets the level of initiation dues.
        \item \paritem{a.ii}{Fund Administration}
        Administers funds available for Tau Beta Pi scholarships according to procedures
established in writing by the officer corps, unless a different method of
administration is specified for the funds.
        \item \paritem{a.iii}{Creation of Ad Hoc Officer Positions}
        May, with Advisory Board approval, create ad hoc officer positions. These positions:
        \begin{enumerate}
            \item Must be listed in Appendix \ref{F.4}.
            \item May be placed on a new or existing Team (see Section \hyperref[I.4.b]{I.4.b}).
            \item May only last for two academic terms, after which time they must be approved
by the chapter membership. This can be either as a permanent officer
position, or as an extension of the ad hoc position. An extension of the
ad hoc position requires a simple majority vote of the chapter membership,
and may not be for longer than another two terms. An officer position may
not exist in an ad hoc state for more than four consecutive terms. Any ad
hoc position that has existed for four consecutive terms cannot be recreated
as ad hoc without a gap of at least two terms.
        \end{enumerate}
        \item \paritem{a.iv}{Creation of Chair Positions}
        May create Chair positions as needed. These positions:
        \begin{enumerate}
            \item Must be listed in Appendix \ref{H}.
            \item May be placed on either the Events Team or Chapter Team (see Section \hyperref[I.4.b]{I.4.b}).
            \item Must report to an Officer or group of Officers, who may be a member(s) of
the Chair’s Officer Team (hereafter referred to as Reporting Officer(s)). Each
Chair’s Reporting Officer(s) shall be listed in Appendix \ref{H}.
        \end{enumerate}
    \end{enumerate}
    \item \subsecitem{b}{Officer Teams}
    The officer corps is organized into teams. These teams are the Executive
Committee, Events Team, and Chapter Team. The lead and membership of
each team is defined in Appendix \ref{F}.
    \begin{enumerate}
        \item \paritem{b.i}{Team Leads}\textsc{Team Leads} 
        Team Leads are responsible for the oversight and coordination
of the team and for ensuring that team members tasks are completed. Team
members retain responsibility for how their tasks should be completed. Team
Leads are additionally responsible for serving as mentors for their team members
when necessary.
        \item \paritem{b.ii}{Purpose}\textsc{Purpose} 
        Teams are intended to assist the President in overseeing the officers,
to create additional opportunities for leadership growth, and to help facilitate
mentoring of newer officers. They are not intended to segment the decision
making of the officer corps, and as such most communication and discussion
should happen amongst the entire officer corps.
        \item \paritem{b.iii}{Executive Committee}\textsc{Executive Committee} 
        In addition to the individual responsibilities of the constituent
officers, the Executive Committee is collectively responsible for setting
the semester calendar and for designating the Convention Delegate (fall term).
    \end{enumerate}
\end{enumerate}


\subsection{Committees} \label{I.5}

\begin{enumerate}
    \item \subsecitem{a}{Purpose}
    Committees are intended to facilitate additional chapter activity and provide
members additional involvement opportunities. Committee function and schedule
will be determined by the chair, with input from the committee members. Committees
are distinct from Teams in that Teams are comprised of officers with similar
but distinct responsibilities, while committees are open to any member, per their selection
procedures, and are, in general, geared toward a more focused task.
    \item \subsecitem{b}{Standing Committees}
    The standing committees are the Professional Development
Committee,Website Committee, K-12 Outreach Committee, and Group Leaders Committee.
The membership and duties of each are listed in Appendix \ref{G.2}.
    \item \subsecitem{c}{Ad Hoc Committees}
    In addition to the standing committees, the officer corps, with
Advisory Board approval, may create ad hoc committees. These committees:
    \begin{enumerate}
        \item \paritem{c.i}{Enumeration}
        Must be listed in Appendix \ref{G.3}.
        \item \paritem{c.ii}{Chair}
        May be chaired by an officer or another active member, as selected by the officer corps.
        \item \paritem{c.iii}{Membership}
        Will be composed of members on a volunteer basis. (In the event of detrimental
participation and at the recommendation of the committee chair, the officer
corps may remove a member of the committee Members so removed may make
an appeal to the Advisory Board, who may reverse the removal by a $\sfrac{5}{7}$ vote.)
        \item \paritem{c.iv}{Duration}May exist for a maximum of two years, after which time they must be approved
by the chapter membership. This can be either as a permanent committee, or
as an extension of the ad hoc committee. An extension of the ad hoc committee
requires a simple majority vote of the body, and may not be for longer than
another year. A committee may not exist as ad hoc for more than three years.
Any ad hoc committee that has existed for three years cannot be recreated as ad
hoc without a gap of at least one year.
    \end{enumerate}
\end{enumerate}


\subsection{Chairs} \label{I.6}

\begin{enumerate}
    \item \subsecitem{a}{Purpose}
    Chair positions are intended to be single-purpose leadership roles within
the chapter that are filled each semester. They allow the leadership of signature
events and established tasks, the addition or subtraction of which would markedly
change the chapter, to be carried out by someone not necessarily an officer.
    \item \subsecitem{b}{Creation}
    Chair positions may be created by the officer corps or the chapter membership
at any time by a simple majority vote at any time. Chair positions may exist
for any length of time, though chairs should be appointed at least semesterly.
    \item \subsecitem{c}{Dissolution and Removal}
    Chair positions may be removed at any time by a $\sfrac{2}{3}$
vote of the officers. Additionally, persons serving as chairs may be removed from
their position by a majority vote of the officers.
\end{enumerate}


\subsection{Terms of Office} \label{I.7}

The officers of this chapter hold office for one semester
except for the External Vice Presidents, one K-12 Outreach Officer, Professional Development
Officer, and Treasurer, whose terms are one calendar year, and the Secretary and
the remaining K-12 Outreach Officer, whose terms are one academic year. Ad hoc officer
positions specified in Appendix \ref{F.4}. will have terms of one semester, unless otherwise
specified. Advisor terms are decided as part of the advisor election procedure, as described
in Bylaw \ref{III.5}. Chair positions specified in Appendix \ref{H} will have terms of one semester,
which concludes at the end of the term of the current Officer Corps.